% Options for packages loaded elsewhere
\PassOptionsToPackage{unicode}{hyperref}
\PassOptionsToPackage{hyphens}{url}
\documentclass[
]{article}
\usepackage{xcolor}
\usepackage{amsmath,amssymb}
\setcounter{secnumdepth}{-\maxdimen} % remove section numbering
\usepackage{iftex}
\ifPDFTeX
  \usepackage[T1]{fontenc}
  \usepackage[utf8]{inputenc}
  \usepackage{textcomp} % provide euro and other symbols
\else % if luatex or xetex
  \usepackage{unicode-math} % this also loads fontspec
  \defaultfontfeatures{Scale=MatchLowercase}
  \defaultfontfeatures[\rmfamily]{Ligatures=TeX,Scale=1}
\fi
\usepackage{lmodern}
\ifPDFTeX\else
  % xetex/luatex font selection
\fi
% Use upquote if available, for straight quotes in verbatim environments
\IfFileExists{upquote.sty}{\usepackage{upquote}}{}
\IfFileExists{microtype.sty}{% use microtype if available
  \usepackage[]{microtype}
  \UseMicrotypeSet[protrusion]{basicmath} % disable protrusion for tt fonts
}{}
\makeatletter
\@ifundefined{KOMAClassName}{% if non-KOMA class
  \IfFileExists{parskip.sty}{%
    \usepackage{parskip}
  }{% else
    \setlength{\parindent}{0pt}
    \setlength{\parskip}{6pt plus 2pt minus 1pt}}
}{% if KOMA class
  \KOMAoptions{parskip=half}}
\makeatother
\setlength{\emergencystretch}{3em} % prevent overfull lines
\providecommand{\tightlist}{%
  \setlength{\itemsep}{0pt}\setlength{\parskip}{0pt}}
\usepackage{bookmark}
\IfFileExists{xurl.sty}{\usepackage{xurl}}{} % add URL line breaks if available
\urlstyle{same}
\hypersetup{
  hidelinks,
  pdfcreator={LaTeX via pandoc}}

\author{}
\date{}

\begin{document}

{
\setcounter{tocdepth}{3}
\tableofcontents
}
Paul Koop

Das Pompeji-Projekt

Vorgeschichte

„Entscheidungen waren noch nicht nötig. Aber der Raum dafür war da.``

Eine szenische Erzählung aus dem Pompeji-Projekt

Wie Michael wurde, was er ist -- und warum Julia ging.

Inhalt

\hyperref[teil-a-boston]{\textbf{TEIL A -- BOSTON 3}}

\begin{quote}
\hyperref[kapitel-a1-letzter-winter-zu-hause]{Kapitel A1 -- Letzter Winter zu Hause 3}

\hyperref[kapitel-a2-der-campus]{Kapitel A2 -- Der Campus 5}

\hyperref[kapitel-a3-die-zweite-sprache]{Kapitel A3 -- Die zweite Sprache 7}

\hyperref[kapitel-a4-eintritt]{Kapitel A4 -- Eintritt 9}
\end{quote}

\hyperref[teil-b-rom-vorluxe4ufige-nuxe4he]{\textbf{TEIL B -- ROM: VORLÄUFIGE NÄHE 11}}

\begin{quote}
\hyperref[kapitel-b1-die-gregoriana]{Kapitel B1 -- Die Gregoriana 11}

\hyperref[kapitel-b2-zusammen-leben]{Kapitel B2 -- Zusammen leben 13}

\hyperref[kapitel-b3-martina]{Kapitel B3 -- Martina 15}

\hyperref[kapitel-b4-die-erste-fahrt-im-vw-bulli]{Kapitel B4 -- Die erste Fahrt im VW-Bulli 16}
\end{quote}

\hyperref[teil-c-der-sanfte-bruch]{\textbf{TEIL C -- DER SANFTE BRUCH 18}}

\begin{quote}
\hyperref[kapitel-c1-die-entscheidung]{Kapitel C1 -- Die Entscheidung 18}

\hyperref[kapitel-c2-abwesenheit]{Kapitel C2 -- Abwesenheit 20}
\end{quote}

\hyperref[teil-d-pompeji]{\textbf{TEIL D -- POMPEJI 21}}

\begin{quote}
\hyperref[kapitel-d1-ankunft]{Kapitel D1 -- Ankunft 21}

\hyperref[kapitel-d2-kindheit-zwischen-ruinen]{Kapitel D2 -- Kindheit zwischen Ruinen 22}

\hyperref[kapitel-d3-archuxe4ologie]{Kapitel D3 -- Archäologie 23}
\end{quote}

\hyperref[teil-e-rom-das-schweigen]{\textbf{TEIL E -- ROM: DAS SCHWEIGEN 24}}

\begin{quote}
\hyperref[kapitel-e1-das-collegium]{Kapitel E1 -- Das Collegium 24}

\hyperref[kapitel-e2-begegnung]{Kapitel E2 -- Begegnung 25}

\hyperref[kapitel-e3-das-schweigen]{Kapitel E3 -- Das Schweigen 26}

\hyperref[kapitel-e4-luca]{Kapitel E4 -- Luca 27}
\end{quote}

\hyperref[teil-f-zwei-kindheiten-parallel]{\textbf{TEIL F -- ZWEI KINDHEITEN (Parallel) 28}}

\begin{quote}
\hyperref[kapitel-f1-martina-luca]{Kapitel F1 -- Martina / Luca 28}

\hyperref[kapitel-f2-sprache]{Kapitel F2 -- Sprache 29}

\hyperref[kapitel-f3-fragment]{Kapitel F3 -- Fragment 30}
\end{quote}

\hyperref[teil-g-der-ethiker-indirekt]{\textbf{TEIL G -- DER ETHIKER (INDIREKT) 31}}

\begin{quote}
\hyperref[kapitel-g1-die-entfernung]{Kapitel G1 -- Die Entfernung 31}

\hyperref[kapitel-g2-ein-gespruxe4ch-am-rand]{Kapitel G2 -- Ein Gespräch am Rand 32}
\end{quote}

\hyperref[teil-h-vor-dem-workshop]{\textbf{TEIL H -- VOR DEM WORKSHOP 33}}

\begin{quote}
\hyperref[kapitel-h1-irarah-entsteht]{Kapitel H1 -- IRARAH entsteht 33}

\hyperref[kapitel-h2-vor-dem-workshop]{Kapitel H2 -- Vor dem Workshop 34}
\end{quote}

\hyperref[personen]{\textbf{Personen 35}}

\begin{quote}
\hyperref[michael-phillips]{Michael Phillips 35}

\hyperref[julia-rossi]{Julia Rossi 35}

\hyperref[martina-rossi]{Martina Rossi 35}

\hyperref[maria]{Maria 36}

\hyperref[luca]{Luca 36}

\hyperref[pater-matteo-conti]{Pater Matteo Conti 36}
\end{quote}

\section{TEIL A -- BOSTON}\label{teil-a-boston}

\subsection{Kapitel A1 -- Letzter Winter zu Hause}\label{kapitel-a1-letzter-winter-zu-hause}

Der Winter kam nicht plötzlich, er schob sich langsam in die Tage. Erst wurde das Licht kürzer, dann die Gespräche. Schließlich lag Schnee auf den Straßen von Boston, festgefahren an den Kreuzungen, grau an den Rändern der Gehwege. Das Haus der Phillips stand ruhig dazwischen, ein schmales Reihenhaus mit Backsteinfassade, nichts Besonderes, nichts Provisorisches.

Michael war siebzehn, vielleicht schon achtzehn, das spielte keine große Rolle mehr. Die Schule lag hinter ihm wie ein bekannter Weg, den man auch im Dunkeln gehen konnte. Morgens stand er früh auf, zog sich an, aß Toast und trank Kaffee, den seine Mutter bereits gekocht hatte. Sie war Biologielehrerin an einer High School im Süden der Stadt, ihr Tag begann strukturiert und endete selten überraschend. Sein Vater, Physiklehrer, verließ das Haus etwas später. Er las die Zeitung gründlich, als müsse er die Welt erst verstehen, bevor er sie verließ.

Am Frühstückstisch wurde wenig gesprochen. Nicht aus Kälte, eher aus Übereinkunft. Man wusste voneinander, ohne es ständig zu sagen. Wenn Michael etwas erzählte, hörten beide Eltern zu. Wenn er schwieg, fragte niemand nach.

An diesem Morgen lag frischer Schnee. Die Straße war noch nicht geräumt, die Autos parkten wie halb versunkene Tiere am Rand. Michael zog den Mantel an, wickelte den Schal zweimal um den Hals.

„Es soll den ganzen Tag so bleiben``, sagte seine Mutter und stellte ihm die Brotdose hin.

„Dann wird es ruhig``, sagte sein Vater, ohne aufzusehen.

Michael nickte. Ruhe war nichts Bedrohliches in diesem Haus. Sie gehörte dazu.

Der Schulweg führte an einer Kirche vorbei, eine von vielen in der Nachbarschaft. Er kannte die Zeiten der Messen, auch wenn er sie nicht auswendig aufsagen konnte. Die Kirche war da, wie der Park oder der Lebensmittelladen an der Ecke. Man ging hin, wenn es passte. Man wusste, dass sie blieb.

Im Unterricht saß Michael meist in der zweiten Reihe. Nicht vorne, nicht hinten. Physik fiel ihm leicht, Biologie auch. In Religion hörte er zu, stellte manchmal Fragen, nie provokativ. Es ging um Verantwortung, um Wissen, um das, was man mit dem tun sollte, was man verstand. Die Antworten blieben vage, aber das störte ihn nicht. Er hatte gelernt, dass nicht alles sofort benannt werden musste.

Nachmittags kehrte er nach Hause zurück, zog die Schuhe aus, ließ den Schnee im Flur schmelzen. Er setzte sich an den Tisch im Wohnzimmer, machte Hausaufgaben, las. Bücher lagen überall, nicht ungeordnet, sondern in kleinen Stapeln. Seine Eltern sammelten sie nicht, sie benutzten sie.

Am Abend aßen sie gemeinsam. Es gab Eintopf, später Tee. Das Radio lief leise.

„Ich habe heute mit den Schülern über Klimamodelle gesprochen``, sagte seine Mutter. „Manche glauben, Zahlen seien Meinungen.``

Sein Vater lächelte kurz. „Das glauben viele.``

„Und du?{\kern0pt}``, fragte sie Michael.

Er überlegte einen Moment. „Zahlen sind Werkzeuge``, sagte er schließlich. „Aber man muss wissen, wofür.``

Sein Vater sah ihn an. „Das ist eine gute Antwort.``

Es lag kein Lob darin, eher Anerkennung.

Später räumten sie gemeinsam auf. Michael trocknete das Geschirr, hörte das Wasser im Spülbecken. Durch das Küchenfenster sah man die Straße, nun leer. Die Schneeflocken fielen dichter, lautlos.

In seinem Zimmer setzte er sich an den Schreibtisch. Er hatte einen Brief von einer Universität in der Schublade, ungelesen. Es war nicht der richtige Moment. Stattdessen schlug er ein Buch auf, las ein paar Seiten, ohne sich alles zu merken.

Er dachte nicht an Weggehen oder Bleiben. Nur daran, dass sich etwas veränderte, ohne dass er es benennen konnte. Wie der Schnee, der fiel und blieb, ohne gefragt zu werden.

Spät am Abend hörte er, wie seine Eltern sich im Wohnzimmer unterhielten, gedämpft, vertraut. Worte über Stundenpläne, über Kollegen, über den nächsten Sonntag. Nichts Dringendes. Gerade deshalb wirkte es verlässlich.

Michael legte sich ins Bett, sah an die Decke. Er spürte keine Unruhe im Sinne von Widerstand. Eher eine Erwartung, die noch keine Form hatte. Etwas, das kommen würde, weil es kommen musste.

Draußen fuhr ein Schneepflug vorbei. Das Geräusch war kurz, hart, dann wieder Stille. Die Straße war nun frei, die Ränder höher. Der Winter hatte sich eingerichtet.

Michael schloss die Augen. Entscheidungen waren noch nicht nötig. Aber der Raum dafür war da.

\subsection{Kapitel A2 -- Der Campus}\label{kapitel-a2-der-campus}

Der Campus lag erhöht, als müsse man ihn ein wenig verlassen, um hierherzukommen. Die Wege zwischen den Gebäuden waren sauber, selbst im Herbst, wenn die Blätter feucht auf dem Asphalt lagen. Michael ging oft zu Fuß, auch wenn der Weg länger war. Es gab ihm Zeit, anzukommen.

Er war achtzehn, im ersten Semester. Theologie und Philosophie standen auf seinem Stundenplan, nebeneinander, ohne Erklärung. Die Vorlesungssäle waren hoch, die Fenster schmal. Wenn der Professor sprach, hallte seine Stimme leicht zurück, nicht störend, eher ordnend. Die Sätze waren klar, die Gedanken langsam. Niemand musste sich beeilen.

Michael saß meist in der Mitte des Raums. Er schrieb mit, nicht alles, nur das, was ihm wichtig erschien. Namen, Begriffe, manchmal ein Satz, der hängen blieb. Neben ihm saßen andere Studenten, manche still, manche eifrig. In den Pausen standen sie zusammen, tranken Kaffee aus Pappbechern, sprachen über Prüfungen, über Schlafmangel, über das Essen in der Mensa.

„Du kommst aus Boston?{\kern0pt}``, fragte ihn einer.

„Ja``, sagte Michael.

„Dann bist du ja schon angekommen.``

Michael lächelte. Es fühlte sich nicht so an, aber er widersprach nicht.

Ein Jesuitenpater leitete das Proseminar am Nachmittag. Er trug einen einfachen schwarzen Pullover, keinen Talar. Seine Stimme war ruhig, fast beiläufig. Er stellte Fragen, ließ Zeit zum Antworten. Michael meldete sich nicht oft, aber wenn er sprach, hörten die anderen zu.

Nach der Stunde blieb er kurz sitzen. Der Pater sammelte die Zettel ein, ordnete sie sorgfältig.

„Sie hören aufmerksam zu``, sagte er, ohne aufzusehen.

Michael wusste nicht, ob es eine Feststellung oder ein Lob war. „Ich versuche es``, sagte er.

Der Pater nickte. „Das reicht fürs Erste.``

Die Bibliothek wurde zu einem festen Ort. Nachmittage dort hatten ein eigenes Tempo. Bücher wurden leise aus den Regalen gezogen, Stühle vorsichtig gerückt. Michael lernte, sich zu verlieren, ohne sich aufzugeben. Er las Texte mehrmals, langsam, manchmal nur wenige Seiten am Stück.

Abends aß er in der Mensa. Das Essen war schlicht. Gespräche verliefen nebenher, selten tief, aber ehrlich. Niemand musste etwas beweisen. Das Studium war Arbeit, keine Bühne.

Manchmal ging Michael über den Campus, wenn es schon dunkel war. Die Gebäude waren beleuchtet, einzelne Fenster noch hell. In der Kapelle brannte oft eine Kerze. Er setzte sich hinein, nicht regelmäßig, nicht aus Pflicht. Er saß einfach da, hörte die Stille.

Es war kein Ort des Zweifels, eher einer des Sammelns. Alles war vorhanden: Wissen, Ordnung, Rituale. Und doch blieb etwas offen. Michael fühlte sich nicht fremd, aber auch nicht angekommen. Er hatte den Eindruck, dass dieser Zustand erlaubt war.

In einem Gespräch mit einem Kommilitonen sagte er einmal: „Ich weiß nicht genau, was ich hier suche.``

Der andere zuckte die Schultern. „Geht mir genauso.``

Das beruhigte ihn mehr als jede Antwort.

Die Semester vergingen. Prüfungen kamen und gingen. Michael bestand sie ohne Aufsehen. Er fiel nicht auf, verschwand aber auch nicht. Er war Teil des Ganzen, ohne sich darin zu verlieren.

Als der erste Winter auf dem Campus kam, lag Schnee auf den Rasenflächen. Studenten liefen schneller, zogen die Schultern hoch. Michael blieb kurz stehen, sah über den Platz. Der Campus war stiller geworden, nicht leer.

Er wusste, dass er hier richtig war. Nicht für immer, vielleicht. Aber für jetzt.

\subsection{Kapitel A3 -- Die zweite Sprache}\label{kapitel-a3-die-zweite-sprache}

Das Gebäude der Physik lag näher an der Straße als die theologischen Fakultäten. Der Eingang war nüchtern, Glas und Beton, kein Ornament. Michael musste einen anderen Weg gehen, eine andere U-Bahn nehmen. Die Fahrt dauerte nicht lang, aber sie markierte einen Übergang.

Er war zwanzig, vielleicht einundzwanzig. Der Stundenplan hatte sich verdichtet. Vormittags Vorlesungen, nachmittags Übungen, dazwischen Kaffee, der bitter schmeckte. Die Räume waren heller als am College, funktionaler. Tafeln voller Gleichungen, die am Ende des Tages wieder weggewischt wurden.

Im Labor trug Michael einen Kittel, der ihm etwas zu groß war. Er stand an einem Tisch mit Geräten, die präzise eingestellt werden mussten. Ein Dozent erklärte den Aufbau, sachlich, ohne Umschweife. Michael hörte zu, stellte Fragen, wenn etwas unklar blieb. Niemand wunderte sich darüber, dass er gleichzeitig Theologie studierte. Es wurde registriert, aber nicht kommentiert.

„Du kommst vom College, oder?{\kern0pt}``, fragte ein Kommilitone, während sie Kabel verbanden.

„Ja.``

„Dann bist du das Denken gewohnt.``

Michael lächelte. „Vielleicht ein anderes.``

Sie arbeiteten schweigend weiter. Zahlen erschienen auf dem Bildschirm, wurden notiert, überprüft. Fehler waren Teil des Ablaufs, nichts Persönliches. Wenn etwas nicht funktionierte, begann man von vorn.

In den Vorlesungen saß Michael weiter vorne. Er wollte sehen, wie die Formeln entstanden. Die Sprache war knapp, präzise. Es gab wenig Raum für Ausschmückung. Ein Satz konnte alles enthalten, wenn er richtig gesetzt war.

Nach der Stunde blieb er manchmal mit anderen Studenten im Flur stehen. Gespräche drehten sich um Aufgaben, um Ergebnisse, um das nächste Praktikum. Es war ein anderes Tempo, schneller, dichter. Michael passte sich an, ohne sich zu verändern.

Abends, zurück in seiner Wohnung, legte er die Hefte nebeneinander. Notizen aus der Theologie, aus der Physik. Er vermischte sie nicht. Es war ihm wichtig, dass jede Sprache ihren eigenen Raum hatte. Er spürte keinen Widerspruch, eher eine Erweiterung.

Ein Dozent blieb eines Tages nach dem Seminar kurz bei ihm stehen. „Sie arbeiten sauber``, sagte er. „Sie lassen sich Zeit.``

Michael nickte. „Ich muss verstehen, was ich tue.``

Der Dozent sah ihn an, prüfend, dann zustimmend. „Das ist selten.``

Auf dem Heimweg ging Michael oft zu Fuß. Die Stadt war laut, selbst abends. Sirenen, Stimmen, Verkehr. Er ließ das Geräusch an sich vorbeiziehen. In seinem Kopf ordneten sich Formeln, Begriffe, Beobachtungen. Es war kein Drängen, eher ein Zusammenfügen.

Manchmal traf er Kommilitonen aus dem College zufällig auf der Straße. Sie grüßten sich, wechselten ein paar Worte. Die Welten berührten sich, ohne zu verschmelzen.

Michael empfand keine Zerrissenheit. Er lernte, dass man nicht alles gleichzeitig aussprechen musste. Manche Dinge verlangten nach Genauigkeit, andere nach Geduld. Beides konnte nebeneinander bestehen.

Als das Semester endete, stand er noch einmal im Labor, sah auf die Geräte, die nun ausgeschaltet waren. Er hatte etwas gelernt, das über die Formeln hinausging. Eine zweite Sprache, die er verstand.

Er verließ das Gebäude, trat hinaus in den Abend. Die Luft war kühl, klar. Boston lag vor ihm, vertraut und offen zugleich.

\subsection{Kapitel A4 -- Eintritt}\label{kapitel-a4-eintritt}

Der Entschluss kam nicht an einem bestimmten Tag. Er hatte keinen Zeitpunkt, an dem man hätte sagen können: Jetzt. Er lag schon eine Weile in Michaels Alltag, wie ein Gegenstand, den man beiseitelegt und immer wieder berührt, um sicherzugehen, dass er noch da ist.

Michael war zweiundzwanzig, vielleicht schon dreiundzwanzig. Er hatte sein Studium geordnet, Abschlüsse geplant, Wege offen. In den Gesprächen mit den Jesuiten wurde nichts gedrängt. Man stellte Fragen, hörte zu, ließ Zeit vergehen.

Ein Pater empfing ihn in einem kleinen Büro am Rand des Campus. Auf dem Schreibtisch lagen keine Aktenstapel, nur ein Notizbuch, ein Stift.

„Sie sind lange geblieben``, sagte der Pater.

Michael nickte. „Ich wollte sicher sein.``

„Sicher wovon?{\kern0pt}``

Michael dachte kurz nach. „Dass ich nicht fliehe.``

Der Pater lächelte leicht. „Das ist eine gute Voraussetzung.``

Es folgten weitere Gespräche. Man sprach über Alltag, über Disziplin, über das Leben in Gemeinschaft. Es ging um Regeln, nicht um Visionen. Michael hörte aufmerksam zu, fragte nach. Er wollte wissen, wie Tage aussahen, wie Jahre vergingen.

Zu Hause erzählte er es seinen Eltern an einem Abend nach dem Essen. Der Tisch war abgeräumt, der Tee noch heiß.

„Ich werde in den Orden eintreten``, sagte er.

Seine Mutter sah ihn an, ruhig. „Wann?{\kern0pt}``

„Im Herbst.``

Sein Vater nickte langsam. „Dann hast du dir Zeit gelassen.``

„Ja.``

Mehr wurde nicht gesagt. Es gab keine Warnungen, keine Erklärungen. Sie wussten, dass Michael keine Entscheidungen traf, um etwas zu beweisen.

Die Vorbereitungen waren unspektakulär. Formulare, Gespräche, Termine. Michael ordnete seine Bücher, gab einige weiter, behielt andere. Er verabschiedete sich von Kommilitonen, ohne große Worte. Manche blieben stehen, fragten nach. Andere nahmen es zur Kenntnis, als wäre es ein Ortswechsel.

Am Tag des Eintritts war es kühl. Ein kleiner Raum, wenige Menschen. Ein schlichtes Ritual, klare Worte. Michael sprach sie ruhig, ohne Zögern. Es fühlte sich nicht wie ein Abschluss an, eher wie ein Übergang.

Nach dem Gottesdienst tranken sie Kaffee. Man sprach über Organisatorisches, über den nächsten Schritt. Rom wurde erwähnt, beiläufig, als etwas, das kommen würde.

Am Abend ging Michael noch einmal durch die Stadt. Die Straßen waren vertraut, die Geräusche dieselben. Er wusste, dass er sie bald verlassen würde. Es machte ihn nicht traurig.

Er stand kurz vor dem Haus seiner Eltern, sah zum Fenster hinauf. Das Licht brannte. Er ging hinein, setzte sich dazu. Es war nichts offen geblieben.

Die Ordnung, für die er sich entschieden hatte, war keine Flucht. Sie war ein Rahmen. Und innerhalb dieses Rahmens blieb er derselbe.

\section{TEIL B -- ROM: VORLÄUFIGE NÄHE}\label{teil-b-rom-vorluxe4ufige-nuxe4he}

\subsection{\texorpdfstring{Kapitel B1 -- \emph{Die Gregoriana}}{Kapitel B1 -- Die Gregoriana}}\label{kapitel-b1-die-gregoriana}

Die Gregoriana lag an diesem Vormittag im Halbschatten. Die Sonne erreichte den Innenhof nur schräg, und die Stimmen der Studierenden verloren sich schnell zwischen den Mauern. Michael war neu in der Stadt, noch unsicher in seinen Wegen. Er kam meist ein paar Minuten zu früh, setzte sich in die zweite Reihe und legte seine Unterlagen ordentlich vor sich hin.

Julia saß zwei Reihen weiter vorn. Sie hatte dunkle Haare, trug sie meist zusammengebunden, nicht aus Absicht, eher aus Gewohnheit. Michael nahm sie zunächst kaum wahr. Erst als der Dozent eine Frage stellte und sie antwortete, hob er den Blick. Ihre Stimme war ruhig, klar, ohne Hast. Sie sprach Italienisch, wechselte dann mühelos ins Englische.

Nach dem Seminar blieben einige Studierende im Raum stehen. Michael sammelte seine Sachen, Julia trat neben ihn.

„Du bist nicht von hier``, sagte sie auf Englisch.

„Boston``, antwortete er.

„Ich dachte mir, dass es nicht Europa ist.``

Sie lächelten beide kurz.

In den nächsten Wochen sahen sie sich regelmäßig. In der Bibliothek, an den Tischen im hinteren Bereich, wo es ruhiger war. Sie arbeiteten nebeneinander, tauschten gelegentlich Bücher. Gespräche entstanden beiläufig, über Texte, über Begriffe, über das Tempo der Vorlesungen.

Einmal saßen sie in einem kleinen Café gegenüber der Universität. Der Kaffee war stark, der Tisch wackelte.

„Die Seminare hier sind langsamer als in den USA``, sagte Michael.

„Dafür reden wir mehr``, antwortete Julia.

„Vielleicht``, sagte er. „Manchmal hilft das.``

Sie lachte leise. „Nicht immer.``

Sie sprachen über ihre Wege nach Rom. Julia war aus Norditalien gekommen, hatte zuvor in Mailand studiert. Michael erzählte vom College, vom Wechsel zwischen Disziplinen. Es klang nicht wie eine Rechtfertigung, eher wie eine Beschreibung.

Sie trafen sich nicht mit Absicht, aber immer wieder. Manchmal gingen sie nach der Bibliothek ein Stück gemeinsam. Sie sprachen über kleine Dinge: das Wetter, den Lärm der Stadt, die Enge der Wohnungen. Nähe entstand nicht aus Gesten, sondern aus Verlässlichkeit.

Eines Abends saßen sie auf den Stufen nahe der Piazza della Pilotta. Die Stadt war laut, der Verkehr dicht. Michael sah auf die vorbeigehenden Menschen.

„Rom ist anstrengend``, sagte er.

Julia nickte. „Aber es bleibt etwas.``

Er wusste, was sie meinte, auch wenn sie es nicht erklärten.

Es gab keine Entscheidung, keinen Moment, der sich absetzte. Sie begannen, ihre Arbeit aufeinander abzustimmen, sich gegenseitig Texte zu schicken, Fragen weiterzugeben. Wenn einer fehlte, merkte es der andere.

Die Beziehung begann nicht als solche. Sie war zuerst Zusammenarbeit. Dann Gewohnheit. Erst später bekam sie einen Namen.

Für den Moment reichte es, nebeneinander zu sitzen und zu arbeiten.

\subsection{\texorpdfstring{Kapitel B2 -- \emph{Zusammen leben}}{Kapitel B2 -- Zusammen leben}}\label{kapitel-b2-zusammen-leben}

Die Wohnung lag im dritten Stock eines Hauses nahe der Via Merulana. Julia lebte dort seit ihrem ersten Semester. Zwei Zimmer, eine schmale Küche, Fenster zum Innenhof. Der Putz bröckelte an manchen Stellen, im Winter zog es durch die Rahmen. Sie hatte sich daran gewöhnt.

Michael kam fast jeden Abend. Er kannte den Weg, die Stufen, den Moment, in dem das Licht im Treppenhaus flackerte. Er trug meist eine Tasche mit Büchern bei sich, manchmal Brot oder Obst vom Laden an der Ecke. Er klingelte nicht lange. Julia öffnete, ließ ihn hinein, stellte keine Fragen.

Sie kochten zusammen. Einfaches Essen, nichts Aufwendiges. Während das Wasser heiß wurde, erzählte einer vom Tag. Von einem Seminar, einer Bemerkung am Rand, einer Statistik, die nicht aufging. Michael sprach ruhig, Julia hörte zu, ergänzte, widersprach gelegentlich.

„Du hast heute wieder zu wenig gegessen``, sagte sie einmal und stellte ihm einen Teller hin.

„Ich habe es vergessen``, antwortete er.

Sie setzten sich an den Tisch, nebeneinander. Die Bücher blieben meist liegen, offen, mit Zetteln dazwischen. Manchmal las einer einen Absatz laut vor. Wenn etwas unklar blieb, ließen sie es stehen.

Michael ging später zurück. Er verabschiedete sich nicht jedes Mal mit Worten. Manchmal nahm er nur seine Tasche, zog den Mantel an. Julia blieb an der Tür stehen, hörte, wie seine Schritte im Treppenhaus verschwanden.

Im Jesuitenhaus war es ruhig. Die Gänge waren lang, die Türen geschlossen. Michael legte die Bücher auf den Schreibtisch, notierte noch ein paar Gedanken. Der Tag war nicht abgeschlossen, aber geordnet.

Am nächsten Morgen sahen sie sich wieder an der Gregoriana. Kurz, beiläufig. Ein Nicken, ein Lächeln. Sie setzten sich nicht nebeneinander, aber sie wussten, wo der andere war.

Am Wochenende blieben sie länger zusammen. Sie gingen durch die Stadt, suchten Wege ohne Eile. Kleine Plätze, schmale Straßen. Sie sprachen über Rom, über das Bleiben und das Weitergehen, ohne etwas festzulegen.

„Wir sind hier für eine Zeit``, sagte Julia einmal.

„Ja``, sagte Michael.

Es war kein Abschied, nur eine Feststellung.

Die Kirche war Teil des Rhythmus. Michael ging zu den festen Zeiten, Julia manchmal mit, manchmal nicht. Es gab keine Erklärungen dafür. Es war kein Thema, das entschieden werden musste.

Abends saßen sie in der Küche, das Fenster offen, Stimmen aus dem Innenhof. Sie sagten wenig. Nähe zeigte sich nicht in Gesten, sondern darin, dass keiner ging, solange der andere noch da war.

Es war kein verborgenes Leben. Auch kein öffentliches. Es war einfach ihr Alltag, zwischen Universität, Stadt und den Wegen, die sie jeden Tag neu gingen.

\subsection{Kapitel B3 -- Martina}\label{kapitel-b3-martina}

Das Krankenhaus lag etwas außerhalb des Zentrums. Der Weg dorthin führte durch Straßen, die Michael kannte, aber an diesem Morgen wirkten sie anders. Der Verkehr war dicht, das Licht grau. Julia saß neben ihm, still, die Hände auf dem Bauch gelegt. Sie sagte wenig. Wenn sie sprach, dann klar.

Im Kreißsaal war es hell. Stimmen, Schritte, kurze Anweisungen. Michael stand an der Seite, hielt Julias Hand, ließ sie los, wenn es nötig war. Er tat, was man ihm sagte, ohne zu fragen. Julia atmete ruhig, konzentriert. Sie wirkte nicht ängstlich, eher gesammelt.

Als Martina geboren wurde, war es kurz still. Dann ein Laut, überraschend kräftig. Julia schloss die Augen, atmete aus. Michael sah das Kind, klein, zusammengerollt, die Haut noch gerötet. Er wusste nicht, was er fühlen sollte. Er wusste nur, dass es da war.

Die Hebamme legte Martina auf Julias Brust. Julia legte den Arm um sie, als hätte sie es schon immer getan. Michael trat näher, sah auf beide hinab. Niemand sagte etwas.

Die ersten Tage waren ungeordnet. Zeiten verschwammen. Michael kam morgens, ging abends, kam wieder. Er lernte, die Bewegungen zu erkennen, das kurze Weinen, das sich wieder legte. Julia ruhte sich aus, stand auf, setzte sich wieder hin. Sie nahm alles an, ohne sich zu beklagen.

In der Wohnung veränderte sich wenig und doch alles. Ein kleines Bett im Schlafzimmer, Decken, Flaschen. Der Tisch blieb derselbe, aber man saß anders daran. Gespräche wurden kürzer, Blicke länger.

Michael hielt Martina zum ersten Mal allein, unsicher. Julia sah ihm zu, korrigierte nicht.

„So ist gut``, sagte sie.

Er nickte, blieb stehen, bewegte sich kaum.

Der Rhythmus stellte sich langsam ein. Schlaf kam in Abschnitten. Essen wurde nebensächlich. Michael passte seine Zeiten an. Er kam früher, blieb manchmal länger. Er ging wieder, wenn es notwendig war.

Draußen ging das Leben weiter. Vorlesungen, Termine, Straßenlärm. Drinnen war etwas entstanden, das keine Erklärung brauchte.

Es gab keine Gespräche über Zukunft oder Ordnung. Nur die Gegenwart. Martina schlief, wachte auf, schlief wieder. Julia beobachtete sie aufmerksam. Michael auch.

Abends saßen sie zusammen, das Licht gedimmt. Martina lag zwischen ihnen, atmete ruhig. Niemand sprach. Es war nichts entschieden worden. Aber etwas war da, das blieb.

\subsection{Kapitel B4 -- Die erste Fahrt im VW-Bulli}\label{kapitel-b4-die-erste-fahrt-im-vw-bulli}

Der Bulli war alt, aber zuverlässig. Hellblau, mit kleinen Roststellen an den Kanten. Michael hatte ihn von einem Kommilitonen übernommen, der Rom verließ. Der Motor klang rau, aber gleichmäßig. Julia mochte dieses Geräusch. Es versprach Bewegung, ohne Eile.

Sie fuhren früh los. Die Stadt lag noch still, die Straßen leerer als sonst. Martina schlief auf dem Rücksitz, in eine Decke gewickelt. Ihr Atem ging ruhig. Julia drehte sich immer wieder zu ihr um, kontrollierte, ob alles in Ordnung war.

„Sie schläft tief``, sagte sie.

Michael nickte, hielt den Blick auf der Straße. „Das mag sie.``

Sie fuhren Richtung Meer, ohne festes Ziel. Die Fenster waren geöffnet, warme Luft kam herein. Michael schaltete das Radio ein, ließ es leise laufen. Musik, die keiner von beiden genau kannte.

An einem kleinen Platz hielten sie an. Bäume spendeten Schatten. Michael parkte den Wagen, stellte den Motor ab. Die Stille danach war fast greifbar.

Sie breiteten eine Decke aus. Julia setzte sich, nahm Martina auf den Arm. Das Kind öffnete kurz die Augen, schloss sie wieder. Michael setzte sich daneben, reichte Julia eine Flasche Wasser.

„Das ist gut hier``, sagte Julia.

„Ja``, sagte Michael.

Sie aßen Brot, Käse, Obst. Nichts Besonderes. Der Bulli stand hinter ihnen, wie ein Schutz. Ab und zu ging jemand vorbei, nickte, sagte nichts.

Am Nachmittag gingen sie ein Stück. Michael trug Martina, vorsichtig, ungeübt. Julia ging neben ihm, hielt seine Schulter kurz fest, wenn der Weg uneben wurde.

„Langsam``, sagte sie leise.

Er lächelte, passte den Schritt an.

Später, als es kühler wurde, fuhren sie weiter. Martina schlief wieder ein. Die Sonne stand tief. Michael fuhr gleichmäßig, ohne zu beschleunigen. Julia lehnte den Kopf ans Fenster, sah hinaus.

Es gab keine Gespräche über das, was kommen würde. Kein Plan, kein Versprechen. Nur der Weg, der sich vor ihnen öffnete.

Am Abend parkten sie an einem einfachen Campingplatz. Michael stellte den Wagen ab, öffnete die Tür. Die Luft roch nach Gras und Staub. Julia legte Martina in den Bulli, zog die Decke zurecht.

Sie saßen noch eine Weile draußen, nebeneinander. Das Licht wurde schwächer. Martina schlief ruhig.

Es war nichts Außergewöhnliches. Aber es war ihres.

\section{TEIL C -- DER SANFTE BRUCH}\label{teil-c-der-sanfte-bruch}

\subsection{\texorpdfstring{Kapitel C1 -- \emph{Die Entscheidung}}{Kapitel C1 -- Die Entscheidung}}\label{kapitel-c1-die-entscheidung}

Sie saßen in der Küche. Es war später Abend, das Fenster stand offen. Geräusche aus dem Innenhof drangen herein, Stimmen, Schritte, ein klapperndes Geschirr. Martina schlief im Nebenraum.

Julia hatte einen Brief vor sich liegen, ungeöffnet. Sie strich mit dem Finger über den Umschlag, ohne ihn anzuheben. Michael saß ihr gegenüber, die Hände um eine Tasse gelegt, die längst kalt war.

„Ich habe heute mit jemandem gesprochen``, sagte Julia.

Michael nickte, wartete.

„In Pompeji. Eine Stelle. In der Beratungsstelle für Familien.``

Sie sah ihn an, ruhig, ohne Erwartung.

„Wann?{\kern0pt}``, fragte er.

„Im Herbst.``

Es folgte eine kurze Stille. Michael sah zum Fenster, dann wieder zu ihr. „Das ist weit``, sagte er schließlich.

„Ja.``

Sie öffnete den Umschlag, zog das Papier heraus, legte es wieder hin. „Es ist eine gute Arbeit. Und ich könnte näher an\ldots`` Sie brach ab, sah in Richtung des Nebenraums.

Michael verstand. Er sagte nichts.

„Ich kann nicht bleiben``, fuhr sie fort. „Nicht so.``

„Ich weiß``, sagte er leise.

Er hatte es gewusst, lange bevor sie es aussprach. Nicht als Gedanke, eher als Gewissheit, die sich langsam eingestellt hatte.

„Ich will nicht, dass es\ldots``, begann sie, hielt inne. „Dass es schwerer wird, als es sein muss.``

Michael nickte. „Ich auch nicht.``

Sie sahen sich an. Es lag keine Bitterkeit zwischen ihnen, eher Müdigkeit. Die Art von Müdigkeit, die entsteht, wenn man etwas lange getragen hat.

„Wir bleiben verbunden``, sagte Julia.

„Ja``, sagte Michael.

„Für Martina.``

„Für sie``, bestätigte er.

Er stand auf, ging ein paar Schritte durch den Raum, blieb stehen. „Ich komme``, sagte er. „So oft ich kann.``

Julia lächelte schwach. „Ich weiß.``

Sie standen noch eine Weile schweigend da. Dann ging Julia ins Schlafzimmer, sah nach Martina. Michael folgte ihr nicht. Er blieb in der Küche, räumte die Tassen weg.

Es war keine Trennung mit Türen, die ins Schloss fielen. Es war ein Schritt, der getan werden musste.

Später, als er ging, nahm er seine Jacke vom Stuhl. Julia stand an der Tür. Sie umarmten sich kurz, fest, ohne Hast.

„Pass auf dich auf``, sagte sie.

„Du auch.``

Die Tür schloss sich leise.

\subsection{Kapitel C2 -- Abwesenheit}\label{kapitel-c2-abwesenheit}

Nach Julias Weggang wurde der Tagesrhythmus klarer. Michael stand früh auf, ging die bekannten Wege, kam abends zurück. Die Wohnung war nicht mehr Teil seines Alltags. Stattdessen verbrachte er mehr Zeit in der Bibliothek, in den Arbeitsräumen der Gregoriana, in stillen Büros mit hohen Regalen.

Er richtete sich einen festen Arbeitsplatz ein. Der Tisch stand nah am Fenster, das Licht fiel gleichmäßig hinein. Bücher lagen offen, Notizen ordentlich gestapelt. Er arbeitete konzentriert, unterbrach sich selten. Wenn er den Raum verließ, ließ er alles so zurück, als würde er gleich wiederkommen.

Die Gespräche mit Kollegen wurden kürzer. Man sprach über Texte, über Termine, über Fristen. Michael war zuverlässig. Er kam vorbereitet, hörte zu, stellte Fragen, wenn es nötig war. Niemand fragte nach seinem Privatleben. Es gab keinen Anlass dazu.

Abends aß er im Refektorium. Die Plätze wechselten, die Gesichter auch. Gespräche verliefen ruhig, oft sachlich. Michael beteiligte sich, ohne sich aufzudrängen. Wenn es still wurde, blieb er sitzen, bis es Zeit war zu gehen.

Manchmal schrieb er Briefe. Kurze Sätze, sachlich. Er fragte nach Martina, nach Julias Arbeit. Die Antworten kamen unregelmäßig, aber sie kamen. Er legte die Briefe in eine Schublade, ordentlich gefaltet.

Die Arbeit nahm zu. Das Projekt mit dem Inventar verlangte Aufmerksamkeit. Rückmeldungen mussten ausgewertet, Anpassungen vorgenommen werden. Michael arbeitete sorgfältig, Schritt für Schritt. Es war keine Flucht. Es war das, was jetzt möglich war.

In der Nacht saß er oft noch am Schreibtisch. Die Gänge waren still, nur vereinzelt Licht hinter Türen. Er schrieb, löschte, schrieb neu. Nicht aus Unruhe, sondern aus Genauigkeit.

Am Wochenende ging er manchmal spazieren. Ohne Ziel. Die Stadt war vertraut geworden, ohne ihm nahe zu sein. Er ging an Cafés vorbei, an Plätzen, an denen er früher gesessen hatte. Er blieb nicht stehen.

Es gab keine Bitterkeit. Nur Abstand. Dinge verschoben sich, langsam, ohne Geräusch.

Michael merkte, dass sich etwas löste. Nicht plötzlich, sondern wie ein Knoten, der sich mit der Zeit von selbst öffnet. Die Verbindung blieb, aber sie trug nicht mehr.

Er arbeitete weiter. Und das genügte.

\section{TEIL D -- POMPEJI}\label{teil-d-pompeji}

\subsection{\texorpdfstring{Kapitel D1 -- \emph{Ankunft}}{Kapitel D1 -- Ankunft}}\label{kapitel-d1-ankunft}

Der Zug hielt ruckartig. Julia nahm die Tasche vom Gepäcknetz, stellte sie kurz ab und hob Martina hoch. Draußen war es warm, trotz der frühen Stunde. Die Luft roch anders als in Rom, dichter, süßer, gemischt mit Staub.

Der Bahnsteig war schmal. Stimmen hallten zwischen den Mauern. Ein Mann rief etwas, das Julia nicht verstand, aber es klang nicht wichtig. Sie ging langsam, achtete auf die Stufen. Martina hielt sich an ihrer Schulter fest und sah sich um, still, wach.

Die Wohnung lag in einer Seitenstraße, nicht weit vom Zentrum. Zwei Zimmer, eine kleine Küche. Die Fensterläden waren geschlossen. Julia öffnete sie, eines nach dem anderen. Licht fiel herein, grell, ungedämpft. Man hörte Autos, Stimmen, irgendwo Musik.

Sie stellte den Koffer ab, setzte Martina auf den Boden. Das Kind ging ein paar Schritte, blieb stehen, setzte sich. Julia lachte leise. Sie nahm einen Lappen aus der Tasche, wischte über den Tisch, öffnete einen Schrank. Es roch nach Feuchtigkeit und Reinigungsmittel.

Am Nachmittag ging sie hinaus. Der Weg zum Supermarkt war kurz. Sie kaufte Brot, Wasser, etwas Obst. Martina saß im Buggy, die Augen halb geschlossen. Die Straße war laut, aber nicht hektisch. Menschen blieben stehen, redeten, gingen weiter.

Abends aßen sie am kleinen Tisch. Julia schnitt das Brot, reichte Martina ein Stück. Draußen wurde es ruhiger. Ein Hund bellte, dann nichts mehr.

Später legte sie Martina ins Bett. Das Zimmer war noch warm. Julia setzte sich auf den Stuhl daneben, blieb eine Weile sitzen. Es war kein Abschied und kein Anfang. Nur ein Ort, der jetzt ihrer war.

Am nächsten Morgen würde sie zur Arbeit gehen. Heute nicht. Heute reichte es, angekommen zu sein.

\subsection{Kapitel D2 -- Kindheit zwischen Ruinen}\label{kapitel-d2-kindheit-zwischen-ruinen}

Martina kannte die Wege zwischen den Ruinen, bevor sie lesen konnte. Julia nahm sie oft mit, wenn sie nach der Arbeit noch einen Umweg machte. Sie gingen langsam, blieben stehen. Martina setzte sich auf niedrige Mauern, strich mit den Fingern über den Stein. Er war warm.

Die Schule lag am Rand der Stadt. Morgens brachte Julia sie hin, nachmittags holte sie sie ab. Martina trug einen Rucksack, der fast zu groß war. Im Klassenzimmer roch es nach Kreide und Staub. Die Lehrerin sprach laut, die Kinder leise durcheinander. Martina hörte zu, meldete sich selten.

Nachmittags spielten sie draußen. Touristen gingen vorbei, hielten Kameras hoch, sprachen fremde Sprachen. Martina achtete nicht auf sie. Für sie waren die Ruinen kein Ziel, sondern Umgebung. Mauern, Schatten, Treppen, die nirgendwohin führten.

Zu Hause gab es einen kleinen Tisch, an dem Julia arbeitete. Akten, Hefte, manchmal ein Buch. Martina saß daneben, malte oder machte Hausaufgaben. Wenn sie fragte, wo ihr Vater sei, antwortete Julia ruhig. Er lebe in Rom, arbeite viel. Mehr sagte sie nicht. Es reichte.

Manchmal kam Post aus Rom. Umschläge, ordentlich beschriftet. Julia legte sie beiseite, öffnete sie später. Martina fragte nicht danach. Der Vater war kein Geheimnis, aber auch kein Teil des Alltags.

An freien Tagen gingen sie früh los. Durch die Straßen, an den geschlossenen Läden vorbei, hinein in das Gelände. Die Stadt war dann still. Martina lief voraus, drehte sich um, wartete. Julia folgte.

Abends saßen sie auf dem Balkon. Geräusche von unten, Stimmen, Geschirr. Martina erzählte von der Schule, von einer Mauer, die sie heute entdeckt hatte. Julia hörte zu. Der Himmel wurde dunkel, langsam.

Es war kein besonderes Leben. Aber es war ihres.

\subsection{Kapitel D3 -- Archäologie}\label{kapitel-d3-archuxe4ologie}

Martina wohnte in einem kleinen Zimmer nahe San Lorenzo. Der Weg zur Universität war kurz, zu Fuß, vorbei an Bars, die schon früh offen hatten. Morgens roch es nach Kaffee und Abgasen. Sie ging meist allein, mit einem Rucksack, in dem Hefte, Bücher und ein altes Maßband lagen.

Die Gebäude der Sapienza wirkten nüchtern. Lange Flure, Treppen aus Stein, Hörsäle mit abgenutzten Sitzen. Martina setzte sich oft in die mittleren Reihen. Sie schrieb mit, langsam, ordentlich. Namen, Jahreszahlen, Schichtenfolgen. Nichts daran war spektakulär. Es war Arbeit.

In den Seminaren sprach sie wenig. Sie hörte zu, stellte Fragen, wenn sie sicher war. Ihre Dozenten kannten sie bald. Nicht als besonders laut, sondern als zuverlässig. Wenn etwas nachzureichen war, tat sie es. Wenn eine Grabung anstand, meldete sie sich.

Die erste Ausgrabung führte sie in die Umgebung von Rom. Früh aufstehen, lange Tage. Erde unter den Fingernägeln, Staub in den Haaren. Sie lernte, Schichten zu unterscheiden, vorsichtig zu arbeiten, nichts zu überstürzen. Abends saßen sie zusammen, aßen einfaches Essen, sprachen über Funde. Martina hörte zu, fragte nach, machte sich Notizen.

Manchmal rief sie ihre Mutter an. Kurze Gespräche. Julia fragte nach dem Studium, nach dem Alltag. Martina erzählte von der Hitze, von den Knochen, von den Mauern, die langsam sichtbar wurden. Über Privates sprachen sie wenig. Es war nicht nötig.

Im zweiten Jahr bekam sie mehr Verantwortung. Sie durfte einen kleinen Bereich selbst betreuen. Sie markierte, dokumentierte, fotografierte. Abends ordnete sie die Daten, verglich Pläne, korrigierte sich. Es gefiel ihr, dass nichts schnell ging.

Zwischendurch fuhr sie nach Pompeji. Sie kannte die Wege, die Ruinen, den Rhythmus der Stadt. Sie ging mit neuen Augen hindurch, blieb stehen, wo sie früher einfach weitergegangen war. Geschichte war jetzt etwas Konkretes.

In Rom, zurück im Zimmer, saß sie oft am Schreibtisch, das Fenster offen. Geräusche von draußen, Stimmen, Motorroller. Sie las, schrieb, strich Sätze. Es gab Tage, an denen sie müde war, und Tage, an denen sie sicher war, dass sie hier richtig war.

Am Ende des Studiums reichte sie ihre Arbeit ein. Kein großer Moment. Sie gab den Ordner ab, ging hinaus, setzte sich auf eine Bank. Die Sonne stand hoch. Sie dachte nicht an das, was kommen würde. Es reichte, dass sie ihren Platz gefunden hatte.

\section{TEIL E -- ROM: DAS SCHWEIGEN}\label{teil-e-rom-das-schweigen}

\subsection{\texorpdfstring{Kapitel E1 -- \emph{Das Collegium}}{Kapitel E1 -- Das Collegium}}\label{kapitel-e1-das-collegium}

\emph{Maria kam früh. Sie schloss die Tür zum Büro auf, stellte ihre Tasche ab und zog die Jalousien hoch. Der Schreibtisch war ordentlich, wie immer. Kalender, Ablagefächer, ein Stapel Briefe.}

\emph{Sie kochte Kaffee in der kleinen Küche am Ende des Flurs. Währenddessen hörte sie Schritte auf den Treppen, Stimmen, gedämpft. Die Seminaristen gingen zum Frühstück. Maria kannte viele beim Namen, wusste, wer pünktlich war und wer nicht. Sie nickte ihnen zu, mehr nicht.}

\emph{Im Büro ordnete sie Post, tippte Termine ab, legte Akten bereit. Deutsch, Ungarisch, Italienisch. Manchmal Englisch. Sie wechselte mühelos. Es war Teil der Arbeit. Die Institution war alt, die Abläufe fest. Maria passte hinein, ohne aufzufallen.}

\emph{Gegen Mittag brachte sie Unterlagen in ein anderes Stockwerk. Sie ging langsam, achtete darauf, niemanden zu stören. Türen standen offen, aus einem Zimmer drang leises Sprechen. Sie blieb nicht stehen.}

\emph{In der Pause saß sie kurz im Hof. Die Sonne hatte inzwischen die Steine erreicht. Sie aß ein Sandwich, sah den Vögeln zu. Der Ort war geschützt, abgeschlossen. Draußen war die Stadt laut, hier nicht.}

\emph{Am Nachmittag kamen Besucher. Priester, Professoren, Gäste aus dem Ausland. Maria empfing sie, nahm Mäntel entgegen, wies den Weg. Sie war freundlich, zurückhaltend. Sie stellte keine Fragen.}

\emph{Abends schloss sie das Büro ab. Sie ging noch einmal durch die Räume, kontrollierte Fenster und Türen. Alles war an seinem Platz. Auf dem Weg nach draußen blieb sie kurz stehen, sah zurück. Das Collegium lag ruhig da, unverändert.}

\emph{Es war nicht ihr Ort, aber sie arbeitete hier. Das genügte.}

\subsection{Kapitel E2 -- Begegnung}\label{kapitel-e2-begegnung}

Sie sahen sich zuerst auf den Fluren. Michael kam aus dem Lesesaal, Maria trug einen Stapel Akten. Sie nickten sich zu, wie man es hier tat. Später tauschten sie ein paar Worte. Termine, ein Formular, eine Unterschrift.

Manchmal begegneten sie sich im Hof. Maria saß auf der Bank, Michael ging vorbei, blieb kurz stehen. Sie sprachen über das Wetter, über den Lärm der Stadt draußen. Nichts Persönliches. Es reichte.

Einmal brachte Maria ihm Unterlagen in sein Zimmer. Bücher lagen auf dem Tisch, Papiere daneben. Sie stellte den Umschlag ab, wartete. Er dankte ihr, fragte, ob sie lange schon hier arbeite. Sie antwortete knapp. Er nickte.

In der Kaffeeküche trafen sie sich häufiger. Michael stand am Fenster, Maria füllte Tassen. Sie tauschten ein paar Sätze. Er fragte nach einer Telefonnummer, sie gab sie ihm. Es ging um eine Rückfrage, sagte sie. Er verstand.

Die Nähe war da, aber sie blieb ruhig. Keine Gesten, keine Worte zu viel. Sie wussten beide, wo sie waren. Die Institution setzte Grenzen, ohne sie auszusprechen.

Als sie sich verabschiedeten, geschah es ohne Anlass. Michael ging schneller, Maria blieb noch einen Moment stehen. Dann ging auch sie weiter.

\subsection{\texorpdfstring{Kapitel E3 -- \emph{Das Schweigen}}{Kapitel E3 -- Das Schweigen}}\label{kapitel-e3-das-schweigen}

Maria merkte die Veränderung früh. Nicht als Schreck, eher als Verschiebung. Sie wurde müder, achtete genauer auf ihren Körper. Im Büro arbeitete sie weiter wie gewohnt, ordnete Akten, beantwortete Anrufe. Niemand fragte.

Sie ging allein zu den Terminen. Die Wartezimmer waren still, sachlich. Sie hörte zu, nickte, unterschrieb. Auf dem Heimweg kaufte sie Brot, Obst, Wasser. Der Alltag blieb.

Michael sah sie weiterhin. Auf den Fluren, im Hof. Er bemerkte, dass sie seltener stehen blieb, dass sie schneller ging. Er fragte nicht. Maria war dankbar dafür.

Als der Bauch sichtbar wurde, wechselte sie ihre Kleidung. Weitere Jacken, weitere Wege. Sie nahm sich früher frei, sagte, es gehe ihr gut. Es stimmte.

Die Geburt war kurz. Ein Raum, weißes Licht, eine Stimme, die ihren Namen sagte. Dann lag das Kind neben ihr, klein, warm. Sie hielt ihn fest, als müsse sie ihn sichern.

Luca schlief viel. Maria sah ihm zu, lernte seine Bewegungen. Sie gab ihm einen Namen, ohne lange zu überlegen. Es war kein Versprechen, nur ein Anfang.

Sie dachte an Michael, manchmal. An seine ruhige Art, an seine Hände, wenn er sprach. Sie dachte auch an das Haus, an die Regeln, an die Wege, die für ihn vorgesehen waren. Sie stellte sich vor, was ein Wort verändern würde. Zu viel.

Also schwieg sie. Nicht aus Angst, nicht aus Trotz. Sie wollte das Kind schützen, und auch ihn. Das Schweigen war eine Entscheidung, kein Mangel.

Als sie mit Luca nach Hause ging, trug sie ihn nah am Körper. Die Stadt war laut, der Himmel klar. Maria ging langsam. Sie hatte Zeit.

\subsection{Kapitel E4 -- Luca}\label{kapitel-e4-luca}

Luca wuchs in einer kleinen Wohnung unweit des Zentrums auf. Die Fenster gingen auf eine schmale Straße, unten parkten Motorroller. Maria stellte sein Bett ans Fenster, damit Luft hereinkam. Er schlief ruhig.

Sie brachte ihn früh in die Betreuung. Morgens gingen sie gemeinsam los, nachmittags holte sie ihn ab. Luca sprach wenig, beobachtete viel. Er setzte sich gern an den Rand, baute mit Bauklötzen, sortierte sie nach Farben. Maria ließ ihn.

Zu Hause gab es Bücher. Alte Ausgaben, Lehrbücher, Kalender. Luca blätterte darin, bevor er lesen konnte. Maria erklärte ihm, was er fragte. Nicht mehr.

In der Schule fiel er nicht auf. Er war weder besonders gut noch besonders schlecht. Die Lehrer sagten, er sei still, zuverlässig. Maria nickte. Sie kannte ihn anders.

Nachmittags blieb Luca oft allein. Maria arbeitete länger, kam erst abends zurück. Luca machte sich etwas zu essen, schaltete den Computer ein. Er lernte, wie Dinge funktionierten. Nicht aus Ehrgeiz, sondern weil es ihn beruhigte.

Sie gingen manchmal ins Collegium. Maria hatte dort noch Kontakt, brachte Unterlagen vorbei. Luca wartete im Hof, sah den Brunnen an, die Mauern. Er kannte den Ort, ohne dazu zu gehören.

Es gab wenige Freunde. Ein Junge aus der Nachbarschaft, später niemand mehr. Luca war nicht unglücklich, aber er war für sich. Er stellte Fragen, die Maria nicht immer beantworten konnte. Dann sagte sie nichts. Auch das war eine Antwort.

Mit sechzehn war er groß geworden, dünn, aufmerksam. Er konnte sich versorgen, sich orientieren. Maria wusste, dass er seinen eigenen Weg gehen würde. Sie begleitete ihn, so lange es ging.

\section{TEIL F -- ZWEI KINDHEITEN (Parallel)}\label{teil-f-zwei-kindheiten-parallel}

\subsection{\texorpdfstring{Kapitel F1 -- \emph{Martina / Luca}}{Kapitel F1 -- Martina / Luca}}\label{kapitel-f1-martina-luca}

Martina kam an den Wochenenden manchmal zurück nach Pompeji. Die Wohnung war dieselbe geblieben. Julia kochte, stellte Teller auf den Tisch, fragte nach dem Studium. Martina erzählte von Seminaren, von einer Grabung außerhalb von Rom, von langen Tagen. Es war ein ruhiges Erzählen, ohne Eile.

Unter der Woche lebte sie in Rom. Die Sapienza hatte ihren eigenen Rhythmus. Morgens füllten sich die Wege, mittags leerten sie sich wieder. Martina saß in der Bibliothek, arbeitete sich durch Texte, machte Skizzen. Sie kannte die Gesichter um sich herum. Man grüßte sich, setzte sich nebeneinander, arbeitete weiter.

Abends ging sie zu Fuß nach Hause. Die Straßen waren laut, vertraut. Sie kaufte Brot, manchmal etwas Käse. Zu Hause schrieb sie weiter, bereitete den nächsten Tag vor. Ihr Leben war dicht, aber überschaubar.

Luca bewegte sich in derselben Stadt, aber anders. Er kannte die festen Zeiten nicht. Er wusste, wann Bibliotheken schlossen, welche offen blieben. Er setzte sich an Tische, an denen niemand Fragen stellte. Er nutzte Rechner, die nicht auf ihn registriert waren.

Tagsüber war er unauffällig. Er ging durch Viertel, in denen man nicht auffiel, wenn man ging. Er hörte Gespräche, merkte sich Namen. Abends wurde die Stadt ruhiger. Luca arbeitete dann. Er schrieb, löschte, begann neu.

Manchmal ging er am Collegium vorbei. Er blieb draußen, setzte sich auf eine Mauer. Drinnen kannte er die Wege, aber sie gehörten ihm nicht. Maria arbeitete dort noch immer. Sie fragten einander wenig.

Martina fuhr mit der Metro, stand dicht gedrängt zwischen anderen Studierenden. Luca vermied sie. Er ging zu Fuß, nahm Umwege. Beide kannten die Stadt, beide bewegten sich sicher. Aber sie nutzten unterschiedliche Türen.

An einem Abend saß Luca in einem Internetcafé. Jemand neben ihm sprach von Mailand, von Kontakten, von einer Fahrt. Luca hörte zu. Er sagte nichts. Rom war noch da, aber nicht mehr genug.

Martina saß zur gleichen Zeit an ihrem Schreibtisch. Sie ordnete Notizen, legte Bücher beiseite. Draußen fuhren Autos vorbei. Sie dachte an den nächsten Tag.

Sie kannten einander nicht. Aber die Stadt hielt beide, für eine Zeit.

\subsection{Kapitel F2 -- Sprache}\label{kapitel-f2-sprache}

Martina lernte früh, genau zu sprechen. In der Schule in Pompeji wurde darauf geachtet. Die Lehrer erwarteten vollständige Sätze, saubere Hefte, pünktliche Abgaben. Julia saß abends am Tisch, half bei Hausaufgaben, las Texte gegen, strich nichts an, stellte nur Fragen. Martina antwortete, manchmal knapp, manchmal ausführlicher. Beides war erlaubt.

Zu Hause wurde Italienisch gesprochen, in der Schule kam Englisch dazu, später Latein. Wörter hatten ihren Platz. Man benutzte sie nicht achtlos. Martina merkte sich Formulierungen, hörte auf Übergänge. Sie mochte es, wenn Dinge richtig klangen.

Als sie nach Rom ging, änderte sich der Ton. An der Universität war Sprache schneller. Dozenten sprachen voraus, setzten Wissen stillschweigend voraus. Martina saß in den Hörsälen, schrieb mit, markierte Stellen, die sie später nacharbeiten musste. In den Seminaren wurde diskutiert. Sie meldete sich nicht oft, aber wenn, dann vorbereitet.

In der Bibliothek herrschte eine andere Ordnung. Bücher wurden zurückgestellt, Arbeitsplätze freigehalten. Martina kannte bald die Etagen, die stillen Tische, die Zeiten, zu denen man besser kam oder ging. Sie legte ihre Materialien nebeneinander, ordnete sie am Ende des Tages wieder ein.

Bei Grabungen lernte sie eine andere Sprache. Kurze Anweisungen, klare Absprachen. Man zeigte auf Fundstellen, maß nach, notierte. Staub lag in der Luft, Stimmen waren gedämpft. Martina arbeitete konzentriert, hörte zu, tat, was nötig war.

Abends schrieb sie Berichte. Sie achtete auf Formulierungen, hielt sich an Vorgaben. Die Universität verlangte Klarheit, keine Ausschmückung. Martina fand sich darin zurecht. Sie wusste, was erwartet wurde.

Institutionen waren für sie keine Fremdkörper. Schule, Universität, Grabungsteams -- sie gaben Struktur. Martina bewegte sich sicher darin. Sie kannte die Regeln, nutzte sie, ohne viel darüber nachzudenken.

Sprache war Teil davon. Sie trug sie durch den Tag, durch Texte, Gespräche, Protokolle. Sie ordnete, verband, hielt fest. Martina sprach nicht viel über sich. Aber was sie sagte, hatte Bestand.

\subsection{Kapitel F3 -- Fragment}\label{kapitel-f3-fragment}

Luca wuchs in Wohnungen auf, die nie ganz fertig wirkten. Mal fehlte ein Regal, mal ein Vorhang. Maria achtete darauf, dass es sauber war. Mehr war nicht immer möglich. Sie arbeitete viel, kam abends müde zurück, stellte Einkäufe ab, fragte nach der Schule. Luca antwortete kurz. Er hatte gelernt, sich selbst zu beschäftigen.

Die Schule wechselte. Nicht oft, aber oft genug, dass Freundschaften dünn blieben. Luca saß meist hinten, beobachtete. Er war kein Störenfried. Er fiel eher dadurch auf, dass er Dinge wusste, die nicht im Stoff standen. Zahlenfolgen, Begriffe, Namen von Programmen. Die Lehrer nahmen es zur Kenntnis, ohne viel nachzufragen.

Nachmittage verbrachte er draußen oder vor Bildschirmen. Internetcafés, später gebrauchte Laptops. Er lernte, wo es Steckdosen gab, wie man Zeit kaufte, ohne viel Geld auszugeben. Er las Foren, Anleitungen, Fragmente von Texten. Er sprang zwischen Themen, blieb selten lange bei einem. Was ihn interessierte, speicherte er ab.

Maria fragte nicht viel. Sie sorgte dafür, dass er morgens aufstand, abends aß. Manchmal saßen sie zusammen, schauten Nachrichten. Luca hörte zu, kommentierte knapp. Politik war für ihn etwas, das im Hintergrund lief, wie Verkehrslärm.

In der Kirche war er selten. Das Collegium kannte er nur von außen. Große Türen, klare Abläufe. Maria sprach nicht darüber. Luca stellte keine Fragen.

Seine Bildung setzte sich aus Stücken zusammen. Ein Onlinekurs hier, ein abgebrochenes Projekt dort. Er lernte schnell, vergaß manches wieder. Er wusste, wie man sich Zugang verschafft, wie man Informationen sortiert. Zeugnisse spielten keine große Rolle.

Nachts war er oft wach. Er schrieb Codeschnipsel, las Chats, folgte Diskussionen, die am nächsten Tag schon wieder verschwunden waren. Kontakte entstanden und lösten sich. Namen blieben vage.

Luca war nicht verloren, aber auch nicht eingebunden. Rom bot ihm Räume, in denen man bleiben konnte, ohne aufzufallen. Er bewegte sich darin sicher. Seine Welt bestand aus Übergängen, offenen Enden, Fragmenten.

\section{TEIL G -- DER ETHIKER (INDIREKT)}\label{teil-g-der-ethiker-indirekt}

\subsection{\texorpdfstring{Kapitel G1 -- \emph{Die Entfernung}}{Kapitel G1 -- Die Entfernung}}\label{kapitel-g1-die-entfernung}

\emph{Die Nachricht kam nicht offiziell. Kein Rundschreiben, keine Versammlung. Michael hörte davon zwischen zwei Terminen, im Flur der Gregoriana, wo Stimmen gedämpft blieben, selbst wenn niemand zuhörte.}

\emph{„Conti ist nicht mehr da``, sagte ein Kollege, fast beiläufig. Er hielt einen Ordner unter dem Arm, als hätte er es eilig. Michael blieb stehen. „Nicht mehr da?{\kern0pt}`` -- „Versetzt. Pensioniert. Irgendetwas in der Art.``}

\emph{In den folgenden Tagen fehlte Conti einfach. Sein Platz im Seminarraum blieb leer. Die Tür zu seinem Büro war geschlossen, das Namensschild noch nicht entfernt. Studenten warteten vergeblich nach der Vorlesung, wie sie es gewohnt waren. Niemand erklärte etwas.}

\emph{Michael erinnerte sich an Gespräche, die nie lange dauerten. Conti sprach leise, stellte Fragen, die offen blieben. Er hatte nie Anweisungen gegeben, nie Positionen verteidigt. Manchmal hatte er nur zugehört, genickt, einen Satz notiert. Michael hatte das als selbstverständlich empfunden.}

\emph{Jetzt änderte sich der Ton. Andere Dozenten übernahmen Themen, sprachen strukturierter, klarer. Es wurde mehr verwiesen, mehr abgesichert. Die Seminare funktionierten weiterhin. Alles funktionierte.}

\emph{Einmal begegnete Michael Conti zufällig in der Nähe der Piazza della Pilotta. Er trug keinen Talar, nur einen Mantel. Sie blieben kurz stehen. Conti fragte nach der Arbeit, nach den Studierenden. Michael antwortete sachlich. Es gab keinen Abschiedssatz.}

\emph{Als sie sich trennten, sagte Conti nur: „Bleiben Sie aufmerksam.`` Nicht eindringlich, nicht mahnend. Wie eine Bemerkung am Rand.}

\emph{Michael ging zurück in die Universität. Der Tag setzte sich fort, als wäre nichts geschehen. Doch etwas hatte sich verschoben. Nicht sichtbar, nicht benennbar. Eher wie ein leiser Verlust an Gewicht, den man erst merkt, wenn man etwas anhebt, das früher schwerer war.}

\subsection{Kapitel G2 -- Ein Gespräch am Rand}\label{kapitel-g2-ein-gespruxe4ch-am-rand}

Der Raum lag im hinteren Teil des Gebäudes. Kein Schild an der Tür, nur ein schmaler Gang, der dort endete. Luca war wegen der Akten gekommen. Jemand hatte ihm gesagt, dass sie hier lagen, ungeordnet, kaum benutzt. Er hatte genickt, nichts weiter gefragt.

Conti saß an einem Tisch nahe dem Fenster. Das Licht war gedämpft, der Nachmittag weit fortgeschritten. Vor ihm lag ein Stapel Papiere, handschriftliche Notizen, nichts Offizielles. Er blickte auf, als Luca eintrat, nicht überrascht, eher aufmerksam.

„Sie suchen etwas?{\kern0pt}``, fragte er.

Luca nannte den Namen, den Zeitraum. Conti hörte zu, stand langsam auf, ging zu einem Regal. Er zog eine Mappe heraus, legte sie auf den Tisch. Seine Bewegungen waren ruhig, ohne Eile.

Luca setzte sich, blätterte. Conti blieb stehen, sah aus dem Fenster. Von draußen drang gedämpfter Lärm herein, Schritte, Stimmen, nichts Bestimmtes.

„Arbeiten Sie hier?{\kern0pt}``, fragte Conti nach einer Weile.

„Nicht direkt``, sagte Luca. Mehr sagte er nicht.

Conti nickte. „Direkt ist selten notwendig``, meinte er, fast beiläufig.

Sie schwiegen. Luca schrieb sich etwas auf, schob die Mappe zurück. Conti setzte sich wieder, faltete die Hände. Er stellte eine weitere Frage, nach dem Studium, nach dem, was Luca interessierte. Luca antwortete knapp, sachlich. Es klang nicht ausweichend, eher vorsichtig.

Conti hörte zu, ohne zu unterbrechen. Schließlich sagte er: „Achten Sie darauf, was Sie nicht einordnen können. Das bleibt oft länger.``

Luca sah kurz auf. Er sagte nichts. Er stand auf, bedankte sich. Conti erwiderte den Gruß mit einem leichten Nicken.

Als Luca den Raum verließ, blieb Conti sitzen. Er nahm seine Notizen wieder auf, schrieb einen Satz zu Ende, als hätte es die Unterbrechung nicht gegeben.

\section{TEIL H -- VOR DEM WORKSHOP}\label{teil-h-vor-dem-workshop}

\subsection{\texorpdfstring{Kapitel H1 -- \emph{IRARAH entsteht}}{Kapitel H1 -- IRARAH entsteht}}\label{kapitel-h1-irarah-entsteht}

\textbf{Szene 1}

Luca saß an einem Tisch in seiner kleinen Wohnung. Die Fenster waren geöffnet, Straßenlärm drang gedämpft herein. Vor ihm lagen Ausdrucke, Notizen, alte Texte aus unterschiedlichen Zusammenhängen. Er las, markierte, legte beiseite. Manche Seiten zerknüllte er und warf sie in den Papierkorb. Andere ordnete er neu. Es gab keine Reihenfolge, die blieb. Er arbeitete langsam, ohne Zielmarke.

\textbf{Szene 2}

Am Bildschirm öffnete er ein Repository. Der Code war nicht von ihm. Er las Kommentare, verfolgte Änderungen. Er fügte einen eigenen Hinweis hinzu, präzise, knapp. Keine Erklärung. Er schloss das Fenster, ohne auf eine Reaktion zu warten.

\textbf{Szene 3}

Im Café an der Ecke saßen drei Personen an einem kleinen Tisch. Tassen klirrten, Stimmen überlagerten sich. Luca sagte wenig. Er stellte eine Frage, hörte die Antworten. Es entstand keine Einigung. Jemand lachte, ein anderer wechselte das Thema. Luca zahlte, stand auf, ging.

\textbf{Szene 4}

Woanders, in einem Projekt, das nichts mit ihm zu tun hatte, tauchte ein Einwand auf. Ein kurzer Absatz, sachlich formuliert. Er wurde gelesen, kommentiert, verschoben. Niemand fragte nach dem Urheber. Die Arbeit ging weiter, leicht verändert.

\textbf{Szene 5}

Später, wieder allein, schrieb Luca einen Namen in eine Datei. Kein Titelblatt, kein Datum. Der Name passte, ohne sich aufzudrängen. Er speicherte, schloss den Rechner. Draußen wurde es dunkel.

\subsection{\texorpdfstring{Kapitel H2 -- \emph{Vor dem Workshop}}{Kapitel H2 -- Vor dem Workshop}}\label{kapitel-h2-vor-dem-workshop}

Michael ging früh durch den Innenhof der Universität. Die Pflastersteine waren noch feucht vom nächtlichen Reinigen. Er trug eine Mappe unter dem Arm, die Schritte gleichmäßig. Im Büro warteten E-Mails, ein Ablaufplan, die Unterlagen für den Workshop. Kollegen nickten ihm zu, kurz, respektvoll. Er setzte sich, öffnete den Rechner, ging die Punkte noch einmal durch. Alles war vorbereitet.

In Pompeji stand Martina am Rand einer Grabungsfläche. Die Sonne war bereits hoch, der Staub lag in der Luft. Sie notierte Maße, gab Anweisungen weiter, überprüfte Funde. Studierende hörten zu, machten ihre Arbeit. In der Pause trank sie Wasser, sah auf die Mauern, die seit Jahrhunderten standen. Sie dachte nicht weiter darüber nach.

In Rom saß Luca in einer Bibliothek, die wenig besucht war. Er wechselte zwischen Dokumenten, Nachrichten, offenen Fenstern. Namen tauchten auf, verschwanden wieder. Er schrieb kurze Nachrichten, erhielt Antworten, speicherte sie ab. Es gab keinen festen Ort, von dem aus er arbeitete, nur Verbindungen.

Michael schloss am Abend die Tür seines Büros. Die Stadt lag ruhig. Er ging nach Hause, las noch ein paar Seiten, legte das Buch weg.

Martina kehrte in ihre Wohnung zurück, wusch den Staub von den Händen, bereitete etwas Einfaches zu essen.

Luca verließ die Bibliothek, mischte sich unter die Menschen auf der Straße.

Alles hatte seinen Platz. Niemand wusste, dass es zusammenhing.

\section{Personen}\label{personen}

\subsection{Michael Phillips}\label{michael-phillips}

Geboren 1973 in Boston, Massachusetts.\\
Aufgewachsen in einem Lehrerhaushalt.\\
Studium der Theologie am Boston College, Physikstudium an der Boston University.\\
Eintritt in den Jesuitenorden, Priesterweihe in Boston.\\
Weiterführende Studien und Lehrtätigkeit an der Pontificia Università Gregoriana in Rom.\\
Vater von Martina und Luca.

\subsection{Julia Rossi}\label{julia-rossi}

Geboren in Italien.\\
Studium der Psychologie an der Pontificia Università Gregoriana in Rom.\\
Lebte während des Studiums mit Michael Phillips in Rom.\\
Mutter von Martina.\\
Arbeitete nach der Trennung als Psychologin in Pompeji.\\
Zog Martina allein groß.

\subsection{Martina Rossi}\label{martina-rossi}

Geboren in Rom.\\
Kindheit und Jugend in Pompeji.\\
Studium der Archäologie in Rom.\\
Arbeitet im archäologischen Kontext.\\
Kennt ihren Vater aus frühen Jahren, wächst ohne ihn auf.

\subsubsection{\texorpdfstring{\hfill\break
}{ }}\label{section}

\subsection{Maria}\label{maria}

Geboren in Italien.\\
Arbeitete als Sekretärin am Pontificium Collegium Germanicum et Hungaricum in Rom.\\
Mutter von Luca.\\
Hielt dessen Herkunft geheim.

\subsection{Luca}\label{luca}

Geboren in Rom.\\
Wuchs bei seiner Mutter auf.\\
Erhielt eine fragmentarische, überwiegend autodidaktische Bildung.\\
Bewegte sich früh in digitalen und informellen Netzwerken.\\
Stand in losem Kontakt zu kirchlichen und akademischen Randorten.

\subsection{Pater Matteo Conti}\label{pater-matteo-conti}

Geboren 1948 in Italien.\\
Jesuit, Moralphilosoph.\\
Lehrte an der Pontificia Università Gregoriana in Rom.\\
Wurde 2005 von seinen Aufgaben dort entbunden.\\
Blieb im kirchlichen Umfeld präsent.\\
Begegnete Luca einmal zufällig.

\end{document}
